\chapter{Ausgaben und Diagnose}\label{ch:ausgabe}

\section{Terminalausgabe von \texttt{fgsolve}}
Für jeden Scanpunkt erscheint ein Block:
\begin{lstlisting}
lambda = 0.6328 um, phi = 10 deg, psi = 0 deg, delta = 90 deg
  side   n    m   efficiency
  R     -1    0   0.2329942258
  R      0    0   0.0901052774
  T     -1   -1   0.0698534457
  ...
  sum R+T = 0.789794795102   absorption = 2.102e-01
\end{lstlisting}
Darunter stehen gegebenenfalls
\begin{itemize}
  \item \texttt{precision fp32 (repeated in fp64)} -- die tatsächlich verwendete Genauigkeit, wenn nicht fp64,
  \item \texttt{warning: \ldots} -- Warnungen (z.\,B.\ \val{FFF\_iso} bei einem Metall),
  \item bei isolierten Strukturen \texttt{scattered up = \ldots{} down = \ldots{} absorbed = \ldots} statt der
        Ordnungen,
  \item bei Modenanregung die Zeilen \texttt{MR}, \texttt{MT} mit effektivem Index und Leistung je Mode,
  \item bei Sweeps vor jedem Wert die Zeile \texttt{== sweep <pointer> = <wert>},
  \item bei dispersiven Materialien je Material und Scanpunkt die ausgewerteten Werte, bei \key{eps\_principal}
        je Hauptachse mit \texttt{[1]}--\texttt{[3]}:
\begin{lstlisting}
  material ag: n = 0.1578809+1.8273612i   eps = -3.3143227+0.57701088i
\end{lstlisting}
\end{itemize}
Am Ende nennt \cmd{fgsolve} die geschriebenen Dateien.

\section{\texttt{efficiencies.csv}}\label{sec:csv}
\index{efficiencies.csv@\file{efficiencies.csv}}%
Eine Zeile je Ordnung und Scanpunkt, Komma als Trennzeichen, Punkt als Dezimalzeichen.
\begin{center}
\small
\begin{tabularx}{\linewidth}{lY}
\toprule
Spalte & Bedeutung \\
\midrule
\texttt{sweep\_parameter}, \texttt{sweep\_value} & nur bei Sweeps: Pointer und Wert \\
\texttt{lambda\_um}, \texttt{phi\_deg}, \texttt{psi\_deg}, \texttt{delta\_deg} & Scanpunkt \\
\texttt{side} & \texttt{R} (reflektiert) oder \texttt{T} (transmittiert); Sonderfälle siehe unten \\
\texttt{n}, \texttt{m} & Ordnung (in 1D ist \texttt{m} $=0$) \\
\texttt{efficiency} & Beugungseffizienz \\
\texttt{eff\_s}, \texttt{eff\_p} & s- und p-Anteil (nur bei isotropen Halbräumen, sonst 0) \\
\texttt{amp\_s\_re}, \texttt{amp\_s\_im}, \texttt{amp\_p\_re}, \texttt{amp\_p\_im} & komplexe Amplituden der
s- und p-Komponente \\
\bottomrule
\end{tabularx}
\end{center}
Sonderwerte von \texttt{side}:
\begin{center}
\small
\begin{tabularx}{\linewidth}{lY}
\toprule
\texttt{side} & Bedeutung \\
\midrule
\texttt{S\_up}, \texttt{S\_down}, \texttt{S\_abs} & isolierte Strukturen: Streuleistung nach oben, nach unten, Absorption
(1D: Streubreite in \textmu m, 2D: Streuquerschnitt in \textmu m$^2$) \\
\texttt{MR}, \texttt{MT}, \texttt{Mrad} & Modenanregung: reflektierte und transmittierte Modenleistung (\texttt{n} =
Modennummer), abgestrahlte Leistung \\
\texttt{C\_R}, \texttt{C\_T}, \texttt{C\_A} & Zylinder, Strahl: Reflexion, Transmission, Rest (bezogen auf den
einfallenden Strahl) \\
\texttt{C\_Rm}, \texttt{C\_Tm} & Zylinder: Anteil je azimutaler Harmonischer (\texttt{n} $=m$) \\
\texttt{C\_Sup}, \texttt{C\_Sdown}, \texttt{C\_Ssca} & Zylinder, ebene Welle: Streuquerschnitte in \textmu m$^2$ \\
\bottomrule
\end{tabularx}
\end{center}

\textbf{Einlesen.}
\begin{lstlisting}[style=py]
import pandas as pd
df = pd.read_csv("out/v3/efficiencies.csv")
t1 = df[(df.side == "T") & (df.n == 1)]
print(t1[["lambda_um", "efficiency"]])
\end{lstlisting}

\section{\texttt{materials.csv}}\label{sec:materials-csv}
\index{materials.csv@\file{materials.csv}}%
Mit \val{"output": \{"materials": "materials.csv"\}} schreibt \cmd{fgsolve} die Materialwerte je
Scan-Wellenlänge, auch die der festen Materialien (ohne \key{eps\_tensor}):
\begin{lstlisting}
lambda_um,material,component,n_re,n_im,eps_re,eps_im
0.4,ag,iso,0.15788090118,1.82736122912,-3.31432268272,0.577010875269
0.4,air,iso,1,0,1,0
\end{lstlisting}
\texttt{component} ist \val{iso} oder bei \key{eps\_principal} die Hauptachse \val{1}--\val{3}. So lässt sich
nachprüfen, mit welchen Werten tatsächlich gerechnet wurde.

\section{Diagnose}\label{sec:diagnose}
\index{Diagnose}%
\begin{center}
\small
\begin{tabularx}{\linewidth}{lY}
\toprule
Größe (Python: \texttt{res.diagnostics[...]}) & Bedeutung, Bewertung \\
\midrule
\texttt{energy\_balance} & $\sum R + \sum T$; bei verlustfreien Strukturen 1 \\
\texttt{max\_condition\_t11} & größte Konditionszahl der Teilblöcke bei der S-Matrix-Verkettung; sehr große Werte
($>10^{12}$) deuten auf numerische Probleme \\
\texttt{rk\_steps} & Zahl der Runge--Kutta-Schritte \\
\texttt{warnings} & Liste von Warnungen (auch: Tabelle außerhalb ihres Bereichs mit \val{"extrapolate":
"constant"}) \\
\texttt{materials} & dispersive Materialien am Scanpunkt: Liste mit \texttt{name}, \texttt{component} (0 isotrop,
1--3 Hauptachse), \texttt{n}, \texttt{eps} \\
\texttt{precision} & tatsächlich verwendete Genauigkeit (\val{fp64}, \val{mixed}, \val{fp32}) \\
\texttt{precision\_fallback} & \val{true}, wenn eine fp32-Rechnung in fp64 wiederholt wurde \\
\texttt{refinement\_steps}, \texttt{refinement\_fallbacks} & Nachiterationen bei \val{mixed} \\
\texttt{timings} & \texttt{solve\_s} und bei \texttt{profile=True} die Zeit je Abschnitt \\
\bottomrule
\end{tabularx}
\end{center}

\section{Profil der Rechenzeit}
\index{Profil der Rechenzeit}%
\begin{lstlisting}[style=shell]
$ fgsolve --profile examples/nearfield_pillars_2d.json
...
profile (busy time summed over threads):
  wall                           1.5766 s         1 x  max 1.5766 s
  nearfield                      1.0279 s         1 x  max 1.0279 s
  nearfield.coefficients         0.8220 s         1 x  max 0.8220 s
  solve.assemble                 0.5543 s         2 x  max 0.2948 s
  solve.layer_models             0.3919 s         2 x  max 0.2352 s
  linalg.eig                     0.1728 s         2 x  max 0.0990 s
  ...
written: out/nearfield_pillars_2d/profile.json
\end{lstlisting}
Die Zeiten sind über alle Threads summiert. Typische Abschnitte: \texttt{linalg.eig} (Eigenwertprobleme der RCWA),
\texttt{smatrix.star} (Verkettung), \texttt{nf.raster}/\texttt{nf.coefficients} (Normalenfeld),
\texttt{nearfield.*} (Nahfeld). In Python: \texttt{fg.solve(cfg, profile=True).diagnostics["timings"]["scopes"]}.

\section{Rückgabewerte von \texttt{fgsolve}}
\index{Rückgabecode}%
\begin{center}
\small
\begin{tabular}{cl}
\toprule
Code & Bedeutung \\
\midrule
0 & Erfolg \\
1 & ohne Argument aufgerufen (Kurzhilfe) \\
2 & Funktion noch nicht umgesetzt (\texttt{NotImplemented}) \\
3 & Fehler in Konfiguration oder Rechnung (\texttt{fgsolve: error: \ldots}) \\
4 & \texttt{-{}-cad-report}: Befunde mit Fehlern \\
\bottomrule
\end{tabular}
\end{center}
In Skripten lässt sich so zuverlässig prüfen, ob eine Rechnung gelaufen ist.
